\documentclass[10pt,a4paper]{article}
\usepackage[utf8]{inputenc}
\usepackage{pgfplots}
\pgfplotsset{compat=1.15}
\usepackage{mathrsfs}
\usetikzlibrary{arrows}
\pagestyle{empty}
\usepackage{amsmath}
\usepackage{amsfonts}
\usepackage{comment}
\usepackage{amssymb}
\usepackage{graphicx}
\usepackage[left=2cm,right=2cm,top=2cm,bottom=2cm]{geometry}
\usepackage{listings}
\usepackage{color}
\definecolor{munsell}{rgb}{0.0, 0.5, 0.69}
\definecolor{dkgreen}{rgb}{0,0.6,0}
\definecolor{gray}{rgb}{0.5,0.5,0.5}
\definecolor{mauve}{rgb}{0.58,0,0.82}
\definecolor{minas}{RGB}{0.244, 0.172, 0.36}
\definecolor{PUJ}{RGB}{44, 86, 151}
\definecolor{PUJ2}{RGB}{43, 93, 156}
\definecolor{PUJ3}{RGB}{20, 52, 107}
\definecolor{cyandk}{rgb}{0.0, 0.72, 0.92}

\author{Iván Andrés Trujillo }

%% code information
\lstset{frame=tb,
  language=Python,
  aboveskip=3mm,
  belowskip=3mm,
  showstringspaces=false,
  columns=flexible,
  basicstyle={\small\ttfamily},
  numbers=none,
  classoffset=1,
  morekeywords={True,False}, keywordstyle=\color{munsell}, 
  classoffset=0, 
  keywordstyle=\color{blue},  
  commentstyle=\color{dkgreen},
  stringstyle=\color{PUJ3},
  numberstyle=\tiny\color{gray},
  breaklines=true,
  breakatwhitespace=true,
  tabsize=4,
}

\begin{document}

\newcommand{\HRule}{\rule{\linewidth}{0.5mm}}

%----------------------------------------------------------------------------------------
%	LOGO SECTION
%----------------------------------------------------------------------------------------

\includegraphics[width = 4cm]{pujshield.eps}\\[0.5cm] 

\begin{center} 
\textbf{\LARGE Introduction to fuzzy logic}\\[0.2cm]
\emph{\LARGE Notes of class}\\[0.3cm] 
\emph{Iván Andrés Trujillo Abella} \\
\textsc{\Large 
}\\[0.2cm] 
\textsc{\large Facultad de ingeniería}\\[0.5cm] 
\HRule \\[0.4cm]
\end{center}
\vspace{1cm}


\section{Fuzzy logic}

linguistic variable of the following elements:

\begin{itemize}
\item $X$ variable name
\item $T(x)$ set of terms 
\item $\varpi$ universe of discourse
\item $G$ syntax rules
\item $M$ semantic rules
\end{itemize}

degree of belonging $x \in A  = \mu_{A}(x) = 1$, otherwise $ x \notin A = \mu_{A}(x) = 0$.
now we can said that $\mu(x) \in [0,1]$ the degree of belonging of $x$ to a any set is a number in the closed interval $[0,1]$.


we can said that $\varpi$ discrete (countable), a set could be written as   $\tilde{A} = \{ \frac{\mu_{\tilde{A}}(x_{1})}{x_{1}} + \frac{ \mu_{\tilde{A}}(x_{2})}{x_{2}}+...+ \frac{\mu_{\tilde{A}}(x_{n})}{x_{n}}   \}$ or in a compact way in $\tilde{A} = \{ \sum_{i}^{n}\frac{\mu_{\tilde{A}}(x_{i})}{x_{i}} \}$.
in continuous $\varpi$ (uncountable) we can write $\tilde{A} = \{ \int \frac{\mu_{\tilde{A}}(x)}{x} \}$.


Until now the major point is the breaking the law of "tercero excluido".

\subsection{Support}

The support of a fuzzy set $S()$, those elements that have a degree of belonging greater than zero  $\mu_{\tilde{A}}(x)>0$.
\begin{equation}
S(A) = \{ x  | \mu_{A}(x) >  0 \} 
\end{equation}
remember that a set is considered regarding a $\varpi$.



\subsubsection{Singleton}
The support is defined only for one value of the all domain for example $x_{0}$ and $\mu_{A}(x_{0})=1$.


\subsection{$\alpha$ cut-off}
suppose that $A$ is fuzzy set of $\varpi$ then by $\alpha \in [0,1]$ $A_{\alpha}$ is:
\begin{equation}
 A_{\alpha}= \begin{cases}
      x \in \varpi | \mu_{A}(x) \geq \alpha, & \text{if}\ \alpha \geq 0 \\
      S(A), & \text{if}\ \alpha = 0
    \end{cases}
\end{equation}






\section{Cartesian product}
by two sets $A$ and $B$ the Cartesian product is
defined as all possible pairs formed with the sets.
\begin{equation}
A X B : \{(a,b) | a \in A \text{ and } b \in B\}
\end{equation}

for instance $A = \{a_{1}, a_{2}  \}$ and $ B =\{b_{1},b_{2},b_{3}\}$ therefore
$A X B = \{(a_{1}, b_{1}), (a_{1},b_{2}), (a_{1},b_{3}), (a_{2},b_{1}), (a_{2},b_{2}), (a_{2},b_{3})\}.$


when you think in two intervals for instance $ A = [2,4]$ and $B=[3,5]$ we can interpreted as  region in $\mathbb{R}^{2}$.




\definecolor{zzttqq}{rgb}{0.6,0.2,0}
\definecolor{ududff}{rgb}{0.30196078431372547,0.30196078431372547,1}
\definecolor{xdxdff}{rgb}{0.49019607843137253,0.49019607843137253,1}
\begin{tikzpicture}[line cap=round,line join=round,>=triangle 45,x=1cm,y=1cm]
\begin{axis}[
x=1cm,y=1cm,
axis lines=middle,
xmin=0,
xmax=5,
ymin=0,
ymax=7,
xtick={-3,-2,...,16},
ytick={-2,-1,...,10},]
\clip(-3.857150822068336,-2.4794396662222686) rectangle (16.07103777419214,10.327812429018422);
\fill[line width=2pt,color=zzttqq,fill=zzttqq,fill opacity=0.10000000149011612] (2,5) -- (4,5) -- (4,3) -- (2,3) -- cycle;
\draw [line width=2pt] (2,0)-- (4,0);
\draw [line width=2pt,color=zzttqq] (2,5)-- (4,5);
\draw [line width=2pt,color=zzttqq] (4,5)-- (4,3);
\draw [line width=2pt,color=zzttqq] (4,3)-- (2,3);
\draw [line width=2pt,color=zzttqq] (2,3)-- (2,5);
\draw [line width=2pt,dotted] (2,3)-- (2,0);
\draw [line width=2pt,dotted] (4,3)-- (4,0);
\draw [line width=2pt,dotted] (2,3)-- (0,3);
\draw [line width=2pt,dotted] (2,5)-- (0,5);
\draw (2.663853299744647,4.302404620463229) node[anchor=north west] {$A X B�$};
\begin{scriptsize}
\draw [fill=xdxdff] (0,3) circle (2.5pt);
\draw [fill=xdxdff] (0,5) circle (2.5pt);
\draw [fill=ududff] (2,0) circle (2.5pt);
\draw [fill=ududff] (4,0) circle (2.5pt);
\end{scriptsize}
\end{axis}
\end{tikzpicture}



\section{Orthogonal extension}
we have two universes of discourses $\varpi_{1}$ and $\varpi_{2}$ we can think in the cartesian product of them $\varpi_{1} X \varpi_{2}$.

\begin{table}[h]
\begin{tabular}{clccc}
$\varpi_{1} / \varpi_{2}$    & $\tilde{A}$  & a_{1}  & a_{2}  & a_{3}            \\
b_{1}                       & 0.3 & (b_{1},a_{1}) & (b_{1},a_{2}) & (b_{1},a_{3}) \\
b_{2}                       &  & (b_{2},a_{1}) & (b_{2},a_{2}) & (b_{2},a_{3}) \\
b_{3}                       & 0.9 & (b_{3},a_{1}) & (b_{3},a_{2}) & (b_{3},a_{3})
\end{tabular}
\end{table}

now the extension is:

\begin{table}[h]
\begin{tabular}{clccc}
$\varpi_{1} / \varpi_{2}$    & $\tilde{A}$  & a_{1}  & a_{2}  & a_{3}            \\
b_{1}                       & 0.3 & 0.3 & 0.3 & 0.3 \\
b_{2}                       & 0 & 0 & 0 & 0 \\
b_{3}                       & 0.9 & 0.9 & 0.9 & 0.9
\end{tabular}
\end{table}


we can establish that $\mu_{\vec{\tilde{A}}}(x,y) = \mu_{\tilde{A}}(x)$

\section{Orthogonal projection}
This problem is inverse to the problem defined previously.


\begin{table}[h]
\begin{tabular}{cccc}
$\varpi_{1} / \varpi_{2}$      & a_{1}  & a_{2}  & a_{3}            \\
b_{1}                        & 0.8 & 0.9 &1 \\
b_{2}                      & 0 & 0 & 0.8 \\
b_{3}                        & 0.4 & 0.5 & 0.7
\end{tabular}
\end{table}

this is only take the major value for each element in the different combinations, for instance for $b_{1} = 1$ and for $a_{1}=0.8$.
We can establish that $\mu_{\overleftarrow{A}}(x) = \max \mu_{A}(x,y)$.



\begin{comment}

outrageous;
nasty:
\end{comment}




\subsection{basic operations}
the Union operation with membership functions.
\begin{table}[h]
\begin{tabular}{lll}
ua & ub & U \\
0  & 0  & 0 \\
0  & 1  & 1 \\
1  & 0  & 1 \\
1  & 1  & 1
\end{tabular}
\end{table}

in summary we can uses $ \mu_{A \cup B}(x) = \max ( \mu_{A}(x), \mu_{B}(x))$.

make the table for intersection and complement.


\subsection{T-norm and S-norm}
remember the unitary cubic




\subsection{fuzzy relations}

We have $A \subset \varpi_{1}$ and $ B \subset \varpi_{2}$ a $R$ is a subset of pairs, the fuzzy relation $\tilde{R} = \frac{\mu_{a_{1}b_{1}}}{(a_{1}b_{1})}+...+$.


in mathematical terms we can said that a fuzzy relation between two fuzzy sets is a subset of the Cartesian product among the universe of discourses. $\tilde{R}  = \tilde{A} X \tilde{B}$ is a subset of the cartesian product of the universe of discourses, and the degree of belonging is defined as $\mu_{\tilde{R}}(x,y) = \min ( \mu_{\tilde{A}}(x), \mu_{\tilde{B}}(y))$.  Remember that we can uses any t-norm.


\subsubsection{Compound relations}
Suppose $\tilde{A}\tilde{R_{1}}\tilde{B}$ and $\tilde{B}\tilde{R_{2}}\tilde{C}$ how we can establish a relation between $\tilde{A}$ and $\tilde{C}$.

\[
\tilde{R_{1}} = \begin{bmatrix}
0.3 & 0.9 & 0.1 \\
0.1 & 0 & 0.4 \\
0.8 & 0.6 & 0.7 
\end{bmatrix}, 
\tilde{R_{2}} = 
\begin{bmatrix}
0.9 & 0.5 & 0.1 \\
0.3 & 0.7 & 0.3\\
0.4 & 0.9 & 0.5
\end{bmatrix}
\]
Think in the composition of two functions $f \circ g(x) = f(g(x))$ therefore $\tilde{R_{1}}\circ \tilde{R_{2}}$ is defined in the a similar way with a multiplication matrix take the $i-th$ row of $\tilde{R_{1}}$ and compare with the $i-th$ column of $\tilde{R_{2}}$   instead of apply the multiplication uses the min operator for instance the $1$ row of $\tilde{R_{1}}$ is $[\tilde{r_{1}}_{11},\tilde{r_{1}}_{12}, \tilde{r_{1}}_{13} ... \tilde{r_{1}}_{1n}]$ and the $1$ column of $\tilde{R_{2}}$ is $[\tilde{r_{2}}_{11},\tilde{r_{2}}_{21}, \tilde{r_{2}}_{31} ... \tilde{r_{2}}_{m1}]^{T}$.
the first element of the composition matrix will be $ \max [ min(\tilde{r_{1}}_{11}, \tilde{r_{2}}_{11}) , min(\tilde{r_{1}}_{12},\tilde{r_{2}}_{21})...min( \tilde{r_{1}}_{1n},\tilde{r_{2}}_{m1}) ]$. applied the steps over all rows and columns respectively.

\[
\tilde{R_{1}} \circ \tilde{R_{2}} = 
\begin{bmatrix}
0.4 & 0.7 & 0.3 \\
0.4 & 0.4 & 0.5 \\
0.8 & 0.7 & 0.5
\end{bmatrix}
\]



\section{Linguistic variables}







\section{fuzzy inference}







\section{Fuzzy C-means}


\end{document}
